%%%%%%%%%%  scan idx 618  |  folio 605  %%%%%%%%%%
\SGAPrintPage{605}%
\[
  R^2f_*(I\otimes f^*M)\simeq \underline{\Hom}_{{O}_S}(Q,M).
\]
\item \uline{The cokernel} \(E\) \uline{of} \(f_*({O}_{X'})\to R^1f_*(I)\) \uline{is a locally free} \({O}_S\)\uline{-module.}
\item \uline{The homomorphism} \((f_*{O}_X)_{\red}\longrightarrow (f_*{O}_{X'})_{\red}\) \uline{is surjective.}
\end{enumerate}
\uline{Then the canonical morphism} \(\underline{\Pic}_{X/S}\longrightarrow \underline{\Pic}_{X'/S}\) \uline{is representable by an affine morphism of finite presentation; more precisely, there is a canonical exact sequence of} f.p.p.f. \uline{sheaves,}
\[
  0\longrightarrow \mathbb V(E^\vee)\longrightarrow \underline{\Pic}_{X/S}
  \longrightarrow\underline{\Pic}_{X'/S}\xrightarrow{u}\mathbb V(Q)
\]
\uline{and, for} \(P=\Ker(u)\)\uline{, the monomorphism} \(P\hookrightarrow\Pic_{X'/S}\) \uline{is representable by a closed immersion and} \(\Pic_{X/S}\) \uline{is a torsor over} \(P\) \uline{under the group} \(V(\check E)\) \uline{(that is, the sequence of} f.p.p.f. \uline{sheaves}
\[
  0\longrightarrow \mathbb V(\check E)\longrightarrow \Pic_{X/S}\longrightarrow P\longrightarrow0
\]
\uline{is exact).}

For an \(S\)-scheme \(T\), consider the following two diagrams with exact rows:
\[
\begin{tikzcd}
0 \arrow[r] & I_T \arrow[r] & {O}_{X_T} \arrow[r,dash] & {O}_{X'_T} \arrow[r] & 0\\
0 \arrow[r] & I_T \arrow[r,"\alpha",dash] \arrow[u,equal] & {O}_{X_T}^{*} \arrow[r] \arrow[u,hook] & {O}_{X'_T}^{*} \arrow[r] \arrow[u,hook] & 0
\end{tikzcd}
\qquad (\alpha(i)=1+i)
\]
\[
\begin{tikzcd}[column sep=tiny]
(f_T)_*({O}_{X'_T}) \arrow[r,"\delta_T"] &
R^1(f_T)_*(I_T) \arrow[d,equal]\\
(f_T)_*({O}_{X'_T}^{*}) \arrow[r,"\delta_T^*"] \arrow[u] &
R^1(f_T)_*(I_T) \arrow[r] &
R^1(f_T)_*({O}_{X_T}^{*}) \arrow[r] &
R^1(f_T)_*({O}_{X'_T}^{*}) \arrow[r] &
R^2(f_T)_*(I_T) .
\end{tikzcd}
\]

Observe carefully that the square on the left of this diagram \uline{is not necessarily commutative}. Nevertheless, we have:
\sgaStableTarget{sga6.e12.para.3-2}\[
  \textbf{3.2.}\qquad \im\delta_T^*=\im\delta_T.
\]


%%%%%%%%%%  scan idx 619  |  folio 606  %%%%%%%%%%
\SGAPrintPage{606}%
Indeed, we may clearly suppose that \(S\) and \(T\) are affine. Take an affine covering \(\{V_i\}\) of \(X_T\). Let \(g\in\Gamma(X_T;{O}_{X'_T}^{*})\), and let \(g_i\) be a lift of \(g|V_i\) in \(\Gamma(V_i;{O}_{X_T}^{*})\). Then \(\delta_T^*(g)\) is the class of the \(1\)-cocycle \(\delta_{ij}^*=(g_ig_j^{-1}-1)\); \(\delta_T(g)\) is the class of \(\delta_{ij}=(g_i-g_j)=g_j\delta_{ij}^*=g\,\delta_{ij}^*\), so we have
\[
  \delta_T(g)=g\,\delta_T^*(g)
\]
taking into account the action of \(\Gamma(X_T,\underline{O}_{X'_T})\) on \(I_T\), hence on \(H^1(X_T,I_T)\).

By ii), we have \(\Gamma(X_T;{O}_{X'_T})=\Gamma(X;{O}_{X'})\otimes_S\Gamma(T,{O}_T)\), and hence \(g=\sum g_\beta\otimes t_\beta\) for \(g_\beta\in\Gamma(X,{O}_{X'})\) and \(t_\beta\in\Gamma(T,{O}_T)\). Now, by (v), \((f_*{O}_X)_{\red}\longrightarrow(f_*{O}_{X'})_{\red}\) is surjective; consequently \(g_\beta=v(h_\beta)+n_\beta\) for \(h_\beta\in\Gamma(X;{O}_X)\), \(v:\Gamma(X,{O}_X)\to
\Gamma(X,{O}_{X'})\) and \(n_\beta\in\operatorname{nilrad}(\Gamma(X,{O}_{X'}))\).

By i), \((n_\beta\otimes t_\beta)I_T=0\). Consequently, \(\delta_T(g)=h\,\delta_T^*(g)\), where \(h=\sum h_\beta\otimes t_\beta\in\Gamma(X_T,\underline{O}_{X_T})\), whence \(\delta_T^*(g)=h^{-1}\delta_T(g)=\delta_T(h^{-1}g)\) and hence \(\im\delta_T^*\subset\im\delta_T\).

Finally, let \(g=\sum g_\beta\otimes t_\beta\in
\Gamma(X_T,{O}_{X'_T})\). Let \(g_\beta=v(h_\beta)+n_\beta\), where \(n_\beta\in\operatorname{nilrad}(\Gamma(X,{O}_{X'}))\); we then have \(\delta_T(g)=\delta_T\bigl(\sum n_\beta\otimes t_\beta\bigr)=\delta_T\bigl(1+\sum n_\beta\otimes t_\beta\bigr)=\delta_T^*\bigl(1+\sum n_\beta\otimes t_\beta\bigr)\), whence \(\im\delta_T\subset\im\delta_T^*\).

By 3.2, iii) and iv), there is an exact sequence functorial in \(T\):
\[
  0\longrightarrow \mathbb V(E^\vee)(T)\longrightarrow
  R^1(f_T)_*({O}_{X_T}^{*})
  \longrightarrow
  R^1(f_T)_*({O}_{X'_T}^{*}) \longrightarrow \mathbb V(Q)(T),
\]
which, on passing to the associated f.p.p.f. sheaves, gives an exact sequence
\[
  0\longrightarrow \mathbb V(E^\vee)\longrightarrow\underline{\Pic}_{X/S}
  \longrightarrow\underline{\Pic}_{X'/S}\xrightarrow{u}\mathbb V(Q).
\]

For \(P=\Ker(u)\), the monomorphism \(P\hookrightarrow\underline{\Pic}_{X'/S}\) is obtained by base change from the unit section \(S\hookrightarrow\mathbb V(Q)\); since \(\mathbb V(Q)\) is representable and separated, it follows that \(P\hookrightarrow\Pic_{X'/S}\) is representable by a


%%%%%%%%%%  scan idx 620  |  folio 607  %%%%%%%%%%
\SGAPrintPage{607}%
closed immersion.

Finally, \(0\longrightarrow \mathbb V(E^\vee)\longrightarrow\underline{\Pic}_{X/S}\xrightarrow{w}P\longrightarrow0\) is a torsor for f.p.p.f.; a fortiori, by f.p.q.c. descent from a ``splitting'' of the torsor, \(w\) is representable by a faithfully flat affine morphism of finite presentation (see: Oort, ``Commutative group schemes'', Proposition (17.4)).

\sgaStableTarget{sga6.e12.cor.3-3}\textbf{Corollary 3.3.} \uline{Under the conditions of} 3.1\uline{, we have}
\begin{enumerate}[label=\roman*)]
\item \uline{If} \(\underline{\Pic}_{(\zar)(X'/S)}\to\underline{\Pic}_{(\fppf)(X'/S)}\) \uline{is an isomorphism, then} \(\Pic_{(\zar)(X/S)}\to\Pic_{(\fppf)(X/S)}\) \uline{is an isomorphism.}
\item \uline{Supposing that the} \(S\)\uline{-scheme} \(T\) \uline{is affine, if} \(\Pic(X'_T)\to\underline{\Pic}_{(\zar)(X'/S)}^{(T)}\) \uline{is surjective, then} \(\Pic(X_T)\to\underline{\Pic}_{(\zar)(X/S)}^{(T)}\) \uline{is surjective.}
\end{enumerate}

Indeed, (i) follows from the following diagram of sheaves for the Zariski topology
\[
\begin{tikzcd}[column sep=small]
0 \arrow[r] \arrow[d,"S"] & \mathbb V(E^\vee) \arrow[r] \arrow[d,"S"] &
  \underline{\Pic}_{(\zar)(X/S)} \arrow[r] \arrow[d] &
  \underline{\Pic}_{(\zar)(X'/S)} \arrow[r] \arrow[d,"S"] &
  \mathbb V(Q) \arrow[d,"S"]\\
0 \arrow[r] & \mathbb V(E^\vee) \arrow[r,dash] &
\underline{\Pic}_{(\fppf)(X/S)} \arrow[r] &
\underline{\Pic}_{(\fppf)(X'/S)} \arrow[r] &
\mathbb V(Q),
\end{tikzcd}
\]
by applying the five lemma to it.

(ii) follows from the following diagram with exact rows
\[
\begin{tikzcd}[column sep=small]
H^1(X_T,I_T) \arrow[r] \arrow[d] &
H^1(X_T,{O}_{X_T}^{*}) \arrow[r,dash] \arrow[d] &
H^1(X_T,{O}_{X'_T}^{*}) \arrow[r,dash] \arrow[d] &
H^2(X_T,I_T) \arrow[d]\\
H^0(T,R^1(f_T)_*(I_T)) \arrow[r] &
H^0(T,R^1(f_T)_*{O}_{X_T}^{*}) \arrow[r] &
H^0(T,R^1(f_T)_*{O}_{X'_T}^{*}) \arrow[r] &
H^0(T,R^2f_*I_T)
\end{tikzcd}
\]
because the outer vertical arrows are isomorphisms, \(T\) being affine.


%%%%%%%%%%  scan idx 621  |  folio 608  %%%%%%%%%%
\SGAPrintPage{608}%
\sgaStableTarget{sga6.e12.cor.3-4}\textbf{Corollary 3.4.} \uline{Let} \(k\) \uline{be a field, and} \(X\) \uline{a proper} \(k\)\uline{-scheme. Suppose that} \(\underline{\Pic}_{X_{\red}/k}\) \uline{possesses a universal invertible sheaf. (For example, by descent theory this holds when every connected component} \(X'_i\) \uline{of} \(X_{\red}\) \uline{has a rational point, for example when} \(k\) \uline{is algebraically closed, or more generally when the} g.c.d. \uline{of the degrees of the} \(0\)\uline{-cycles on every} \(X'_i\) \uline{is} \(1\)\uline{). Then, over every affine open subset} \(V\) \uline{of} \(\Pic_{X/k}\)\uline{, there is a universal invertible sheaf on} \(X\times_kV\).

Finally, we prove I in the following more precise form:

\sgaStableTarget{sga6.e12.prop.3-5}\textbf{Proposition 3.5.} \uline{Let} \(S\) \uline{be an integral noetherian scheme, and} \(f:X\hookrightarrow Y\) \uline{a nilpotent immersion of proper} \(S\)\uline{-schemes. Then there is a nonempty open subset} \(V\) \uline{of} \(S\) \uline{such that the morphism}
\[
  f^*:\Pic_{Y/S}|V\longrightarrow\Pic_{X/S}|V
\]
\uline{is representable by an affine morphism of finite presentation; in fact,} \(N=\Ker f^*\) \uline{is an affine group scheme with unipotent fibers (in the sense of} SGA 3 XVII).

Indeed, \(X\) is defined by a nilpotent ideal \(I\). Then \(I\) has a filtration by ideals \(I_0=I\supset I_1\supset\ldots\supset I_n=0\) such that \(NI_{i-1}\subset I_i\), where \(N\) is the nilradical of \({O}_Y\) (for example \(I_k=N^kI\)). Let \(X_i\) be the subscheme of \(Y\) defined by \(I_i\). By induction it suffices to study the immersions \(X_{i-1}\hookrightarrow X_i\). We may therefore suppose that \(N.I=0\); this is condition i) of 3.1, and we finish by reducing to the case where the other conditions (ii) to (v) also hold. By the generic flatness theorem (EGA IV 6.9.1), applied to \(\underline{O}_{X'}\), to the \(R^if_*(I)\) and to \(E\), and taking account of the Künneth relations EGA III 6.9.9, we may suppose, after shrinking \(S\), that conditions ii), iii) and iv) hold.


%%%%%%%%%%  scan idx 622  |  folio 609  %%%%%%%%%%
\SGAPrintPage{609}%
There is a finite extension \(K\) of \(k(S)\) such that \((f_K)_*({O}_{X'_K})\) has all its residue field extensions trivial. For an integral scheme \(S'\) finite and flat over an open subset \(V\) of \(S\), with function field \(k(S')=K\), there is an open subset \(V'\) of \(S'\) such that
\[
  \bigl[(f_{S'})_*({O}_{X'_{S'}})\bigr]_{\red}
  \simeq
  \bigl[(f_{S'})_*({O}_{X_{S'}})\bigr]_{\red}
\]
over \(V'\); after shrinking \(S\) and \(S'\), we may suppose that \(S=V\) and \(S'=V'\). Since the assertion to be proved is local for f.p.q.c., we may replace \(S\) by \(S'\). Finally, apply 3.1.

\sgaStableTarget{sga6.e12.rem.3-6}\textbf{3.6. Remark.} It follows in the same way, taking into account the exact sequences \(\underline{\Pic}_{X_i/S}\longrightarrow\underline{\Pic}_{X_{i-1}/S}\longrightarrow\mathbb V(Q_i)\), that, for every \(s\in S\),
\[
  \operatorname{coker}\bigl(\underline{\Pic}^0_{Y_s/s}\longrightarrow\underline{\Pic}^0_{X_s/s}\bigr)
\]
is a unipotent group.

\sgaStableTarget{sga6.e12.sec.4}\section*{4. Proof of II*: the most delicate part of the proof}

Recall assertion II*: Theorem 1.1 when \(X\) and \(Y\) are integral over \(S\) and the sequence (*) \({O}_Y\longrightarrow f_*{O}_X\rightrightarrows g_*({O}_{X\times_YX})\) is exact; applying the generic flatness theorem EGA IV 6.9.1 and EGA III 7.8.8 reduces us to proving the

\sgaStableTarget{sga6.e12.prop.4-1}\textbf{Proposition 4.1.} \uline{Let} \(S\) \uline{be an integral noetherian scheme,} \(\lambda:X\to S\) \uline{and} \(\mu:Y\to S\) \uline{proper morphisms of finite presentation,} \(f:X\to Y\) \uline{a surjective} \(S\)\uline{-morphism, and} \(g:X\times_YX\to Y\) \uline{the canonical morphism. Suppose that}
\begin{enumerate}[label=\arabic*)]
\item \(X,Y,X\times_YX\) \uline{are flat over} \(S\).
\item \(\lambda_*{O}_X={O}_S\), \(\mu_*{O}_Y={O}_S\) \uline{universally.}

